\section{Complete universal coefficients}\label{sec:calculation}

\subsection{The source genus-zero input}
We first calculate connected genus-zero invariants on $Y_+$ without
using the smoothing answer. The charge matrix and triple intersections
are printed in Appendix~\ref{app:data}. The relevant mirror theorem
is \cite[Theorem 25 and Corollary 28]{CCIT}. Removing ray $10$
gives a projective simplicial 24-cone fan and a smooth canonical
toric Deligne--Mumford stack. Its nonunimodular maximal cones are
$1369,1389,1459,1489,1569$, all of index two. Every nonsmooth
orbit has a singleton Newton face and misses the chosen hypersurface.
The anticanonical support function is integral and convex: the line
bundle is nef, Cartier on the coarse space, and pulled back to the
stack. On a genus-zero orbicurve it has trivial stabilizer action and
nonnegative degree on each component, hence vanishing first cohomology.
It is therefore convex in the cited theorem. Its transverse zero
locus is the ordinary smooth threefold $Y_+$, all twisted-sector
classes restrict to zero, and the Calabi--Yau asymptotics give
$J=I/W_0$. The 29-cone refinement changes only the missed locus.
The charge lattice is saturated and the ambient Mori cone is the
nonnegative orthant in these coordinates. These checks justify
using the untwisted Gamma series for the actual source invariants.

For the charge matrix $Q$ put $Q_{a*}=\sum_iQ_{ai}$ and
\[
 c(n+\rho)=\frac{\Gamma(1+\sum_a(n_a+\rho_a)Q_{a*})}
 {\prod_i\Gamma(1+\sum_a(n_a+\rho_a)Q_{ai})}.
\]
Let $W_0,C_a,\widehat C_{ab}$ be the sums of the scalar, first and
second derivatives at $\rho=0$ of $c(n+\rho)/c(\rho)$ times $z^n$.
This quotient removes the universal zeta constants from the derivatives.
Let $a_n$ denote the coefficients of the genus-zero dilogarithm
expansion, indexed by the curve class $n$.
The mirror coordinates and prepotential equations are
\begin{equation}\label{eq:mirror}
 q_a=z_a\exp(C_a/W_0),\qquad
 \sum_{n>0}n_a a_n\operatorname{Li}_2(q^n)
 =\frac12\sum_{b,c}\kappa_{abc}
 \left(\frac{\widehat C_{bc}}{W_0}-\frac{C_bC_c}{W_0^2}\right).
\end{equation}
These are the Frobenius formulas in \cite[Section 4.4 and Appendix A]{DMM}.
Derivatives at reciprocal-Gamma zeros are retained.

All calculations are exact in
\[
 \QQ[z_1,\ldots,z_6]/(z_1^3,z_2^2,z_3^2,z_4,z_5^6,z_6^6).
\]
Its 432 monomials contain all the candidates in the three relevant
bidegrees and every componentwise smaller exponent. Products,
inversion and mirror substitution cannot bring an omitted exponent
back into this set. Solving all six equations \eqref{eq:mirror}
and extracting the dilogarithm coefficients gives the following
complete nonzero table, where $(a,b,c)=(n_1,n_5,n_6)$:
\begin{equation}\label{eq:gwtable}
\begin{array}{c|rrr}
(a,b,c)&(n_2,n_3)=(1,0)&(0,1)&(1,1)\\ \hline
(0,0,0)&1&-2&3\\
(0,0,1)&-2&1&-2\\
(1,1,1)&-2&1&-2\\
(1,1,2)&3&0&4.
\end{array}
\end{equation}
All other entries of these three sectors in the finite quotient
vanish. Every class in the table is indivisible, so $a_n=N^Y_{0,n}$.
The finite support bounds ensure that this calculation covers every
normalized contribution in those sectors.

\subsection{From the numerical table to universal coefficients}
Put $A=\Z_{Y_+}/\cE_{Y_+}$ and $L=\log A$, and work on the face
$n_4=0$. Normalization makes the entire bidegree-$(0,0)$ sector
equal to one. Each nonempty exceptional local coefficient has PT
realization in $z\QQ[[z]]$ by Lemma~\ref{lem:multiple}; the same is
true of each nonconstant coefficient of its inverse series, since
the inverse uses finite products at fixed source degree.
Theorem~\ref{thm:euler} thus puts every coefficient of $A$ in the
three stated bidegrees in $z\QQ[[z]]$ after PT evaluation.
The same holds for $L$: its only additional lower-bidegree products
have bidegrees $(1,0)$ and $(0,1)$ and start at $z^2$.

Each $L_\gamma$ is primitive and, by indivisibility, killed by
$\rho_d$ for $d>1$. Lemma~\ref{lem:pt} therefore implies
\begin{equation}\label{eq:logcoeff}
 L_\gamma=N^{Y_+}_{0,\gamma}x_{0,1}
\end{equation}
in all three sectors. This is the point at which the arbitrary
exceptional multiplicities in Theorem~\ref{thm:euler} are essential.
The numerical table alone would not exclude positive-genus primitive
generators. Higher-genus GW evaluations of $x_{0,1}$ can still be
nonzero.

Set $x=x_{0,1}$, $r=Q^{b_2}$, $s=Q^{b_3}$,
$u=Q^{b_1+b_5}$, $v=Q^{b_6}$. Equations
\eqref{eq:gwtable}--\eqref{eq:logcoeff} give, through these bidegrees,
\[
 L=x(rP_{10}+sP_{01}+rsP_{11})\pmod{(r^2,s^2)},
\]
where
\[
 P_{10}=1-2v-2uv+3uv^2,\quad P_{01}=-2+v+uv,\quad
 P_{11}=3-2v-2uv+4uv^2.
\]
Thus $[rs]A=xP_{11}+x^2P_{10}P_{01}$. Exactly eight of the 37
candidates have nonzero coefficients:
\begin{equation}\label{eq:universal-table}
\begin{array}{c|l}
 \gamma&(A)_\gamma\\ \hline
 (0,1,1,0,0,0)&3x-2x^2\\
 (0,1,1,0,0,1)&-2x+5x^2\\
 (0,1,1,0,0,2)&-2x^2\\
 (1,1,1,0,1,1)&-2x+5x^2\\
 (1,1,1,0,1,2)&4x-10x^2\\
 (1,1,1,0,1,3)&3x^2\\
 (2,1,1,0,2,2)&-2x^2\\
 (2,1,1,0,2,3)&3x^2.
\end{array}
\end{equation}
The remaining 29 coefficients vanish. Four of the nonzero universal
coefficients have zero connected genus-zero value; they arise from
the lower-bidegree product. Since $P_{10}(1,1)=P_{01}(1,1)=0$ and
$P_{11}(1,1)=3$, finite specialization gives
\begin{equation}\label{eq:sum}
 \sum_{\pi\gamma=\beta}(A)_\gamma=3x_{0,1}.
\end{equation}
The common-blow-up comparison transfers the same normalized table
to $Y_-$. No smoothing invariant was used in its derivation.

\subsection{Independent calculation on the plane}
\begin{proposition}\label{prop:target}
For the smoothing plane of Proposition~\ref{prop:plane},
\[
 U_{S,\beta}=3x_{0,1},\qquad
 U_{S,2\beta}=3x_{0,2}+3x_{0,1}^2-6x_{0,1}.
\]
\end{proposition}
\begin{proof}
All cycles in these classes lie on $D$. In degree one, purity
forces the curve scheme to be a reduced line. In degree two a
generically reduced curve is a plane conic. A pure curve $Z$ with
cycle $2L$ satisfies $I_L^2\subset I_Z\subset I_L$: the image of
$I_L^2$ vanishes generically and hence vanishes by purity.
The rank-one torsion-free sheaf $I_L/I_Z$ is $\mathcal O_L(a)$.
The normal sequence splits as
$N_{L/S}=\mathcal O_L(1)\oplus\mathcal O_L(-3)$, so there is a
surjection $\mathcal O_L(-1)\oplus\mathcal O_L(3)\to\mathcal O_L(a)$.
It forces $a\ge-1$, whence $\chi(\mathcal O_Z)=a+2\ge1$.
At equality the $\mathcal O(3)$ subbundle maps to zero, so the
double line is scheme-theoretically in $D$. Thus every stable pair
in degrees one or two has Euler characteristic at least one, and
at equality is a line or conic with its structure sheaf.

Scheme-theoretically the minimal pair spaces are
$P_1(S,\beta)=(\PP^2)^*$ and $P_1(S,2\beta)=\PP^5$.
Indeed surjectivity of the pair section follows by Nakayama, giving
the Hilbert functor. For a degree-$k$ plane divisor $C$, $k=1,2$,
including a double line, the normal sequence is
\[
 0\to\mathcal O_C(k)\to N_{C/S}\to\mathcal O_C(-3)\to0.
\]
The last sheaf has no sections and $h^0(\mathcal O_C(k))$ is $2$
or $5$. Thus the ambient Hilbert scheme has exactly the tangent
dimension of its smooth closed plane-curve subscheme, which exhausts
its points; its local rings are regular and the immersion is an
isomorphism. The symmetric obstruction theory gives counts
$(-1)^2e(\PP^2)=3$ and $(-1)^5e(\PP^5)=-6$
\cite[p.~3]{PT}.

The primitive extremal class $\beta$ has no nontrivial effective
splitting. Compatibility and Lemma~\ref{lem:pt} give
$U_{S,\beta}=3x_{0,1}$. In degree two, compatibility makes
$R=U_{S,2\beta}-3x_{0,2}-3x_{0,1}^2$ primitive and killed by
every $\rho_d$, $d>1$: the removed coproduct cross terms total
$(3+6)x_{0,1}\otimes x_{0,1}$. Lemma~\ref{lem:multiple} gives
$\PT(x_{0,2})\in z^3\QQ[[z]]$, while
$\PT(x_{0,1}^2)\in z^2\QQ[[z]]$. Consequently
$\PT(R)=-6z+O(z^2)$ with no nonpositive powers, and
Lemma~\ref{lem:pt} gives $R=-6x_{0,1}$.
\end{proof}

This verifies \eqref{eq:sum} using different enumerative input. The
two target coefficients are full universal polynomials, not genus
truncations. The source coefficients over $2\beta$ are not calculated
here; the 183-class bound alone makes no assertion about their values.
